\section{Literature Review}
\label{sec:litreview}

This section reviews the literature that situates the thesis and develops, in detail, the two gaps introduced earlier. It moves from the scale of the problem and the carbon signals that make scheduling possible, through the temporal and spatial flexibility the thesis exploits, to the contrast between heuristic and optimization-based methods, and closes by stating precisely how the present work is positioned.

\subsection{The Carbon Footprint of Computing and AI}
\label{sec:lr-footprint}

The motivation for carbon-aware scheduling begins with the scale of the problem. Data centers account for an estimated one percent of global electricity use, a figure that held stable through the 2010s as efficiency gains offset rising demand \citep{masanet2020recalibrating}. The rapid expansion of artificial intelligence is now raising it, with dedicated AI compute a fast-growing contributor whose projected demand is substantial within a few years \citep{devries2023growing}. Quantifying this footprint at the workload level has become an active research area, progressing from early measurements of training emissions \citep{strubell2019energy} to standardized accounting methods \citep{lannelongue2021greenalgorithms, henderson2020towards} and end-to-end models for large language models \citep{faiz2024llmcarbon, luccioni2023estimating}.

Across these accounting studies, a consistent theme emerges that motivates the rest of this thesis. Operational emissions depend not only on how much computation is performed but also on when and where the electricity is drawn, so that the same job run in a different region or at a different hour can emit several times more or less carbon \citep{dodge2022measuring, patterson2022carbon}. This locational and temporal sensitivity is what carbon-aware scheduling exploits. Operational carbon is not the only relevant dimension, however. A parallel literature stresses that the embodied emissions of hardware manufacturing \citep{gupta2021chasing} and the water consumed for cooling and electricity generation \citep{li2023thirsty} are also material, and although both lie outside the scope of this thesis, they motivate its limitations and future work. Together these strands sit within the broader agenda of aligning artificial intelligence with climate-change mitigation \citep{kaack2022aligning}.

\subsection{Carbon-Intensity Signals}
\label{sec:lr-signals}

Exploiting the when-and-where sensitivity identified above first requires a signal that quantifies how clean electricity is at a given time and place. The most common such signal is grid carbon intensity (CI), the average emissions per unit of electricity consumed, which is published at hourly granularity and used directly as a scheduling input in production systems \citep{radovanovic2023}. Because CI is an average measure, a more marginal alternative has been proposed in the form of locational marginal emissions, which attribute emissions to the specific generator that responds to an incremental load and can differ materially from average CI in large grids \citep{cote2024locational, lindberg2022geographic}. A different signal again arises from the 24/7 carbon-free energy (CFE) accounting movement, which matches each hour of consumption against carbon-free generation in the same hour rather than against annual renewable certificates \citep{riepin2024cfe247costs}. This thesis draws on both CI and the carbon-free energy fraction, using the latter to distinguish electricity that is clean because of genuine renewable supply from electricity that is only incidentally low-carbon.

Because schedules must be planned before the relevant hour arrives, a supporting literature forecasts carbon intensity using neural, ensemble, and uncertainty-aware methods \citep{maji2022carboncast, yan2025ensembleci, li2024uncertainty}. The present work treats CI and CFE as given inputs and isolates the scheduling problem, so forecasting is complementary to, rather than part of, the optimization model developed here.

\subsection{Temporal Workload Shifting}
\label{sec:lr-temporal}

With a carbon signal in hand, the first way to act on it is temporal, deferring delay-tolerant work to hours of lower carbon intensity. This lever was characterized in detail by \citet{wiesner2021letswait}, who quantified achievable savings for delay-tolerant workloads across several national grids and found them to depend strongly on the region and on the available delay budget. The most prominent production realization is Google's carbon-aware computing platform, which shifts flexible batch load across the day using Virtual Capacity Curves as a scheduling interface and reports meaningful operational-carbon reductions at scale without violating service objectives \citep{radovanovic2023}. Building on this idea, subsequent systems work has operationalized temporal shifting through cluster schedulers for machine-learning workloads \citep{xu2025green}, elastic resource scaling that varies capacity across carbon windows \citep{hanafy2025}, and system-level facilities that cap or shape application-level emissions \citep{thiede2023carboncontainers, souza2023ecovisor}.

Despite these advances, a recurring finding qualifies temporal shifting and motivates the present thesis. Several studies show that shifting in time alone yields limited and context-dependent savings, because the benefit shrinks once the delay budget exceeds a few hours and is small in grids whose carbon intensity varies little within a day \citep{sukprasert2024limitations, wiesner2021letswait}.

\subsection{Spatial Workload Routing}
\label{sec:lr-spatial}

Because temporal flexibility alone is bounded, attention turns to the spatial lever, which routes flexible work to cleaner regions. Its foundations lie in the geographical load-balancing literature, which established that distributing load across locations according to local grid conditions can reduce both cost and emissions \citep{liu2012, lindberg2022geographic}. More recent carbon-aware variants route requests to data centers with cleaner electricity subject to latency limits \citep{souza2023casper}, extend the idea to federated and geo-distributed machine learning \citep{bian2024cafe, li2025llm}, and add fairness so that no single region bears a disproportionate environmental burden \citep{li2024equitable}.

A closely related strand goes further and optimizes the spatial and temporal dimensions jointly, which is the strand this thesis contributes to. The most directly comparable work is that of \citet{riepin2024spatiotemporal}, who formulate a linear program for spatio-temporal load shifting across European data centers; their objective, however, is 24/7 CFE matching through dedicated procurement rather than dispatch against real-time grid carbon intensity, and they report an optimal data-center separation of three to four hundred kilometres. Other joint formulations adopt mixed-integer models to add multi-resource footprints \citep{attenni2024shifting} or precedence-constrained workloads \citep{lechowicz2025pcaps}. Even so, a recent review concludes that genuine joint spatio-temporal optimization remains comparatively rare and that most methods still address one dimension at a time \citep{asadov2025review}. The present thesis differs from the joint work that does exist in two respects, in that it uses real multi-year grid CI and CFE signals as the dispatch basis and decomposes the realized saving into its spatial and temporal parts.

\subsection{Heuristic versus Optimization-Based Scheduling}
\label{sec:lr-methods}

Whether they exploit temporal, spatial, or joint flexibility, the approaches surveyed so far divide along a methodological line that constitutes the second gap this thesis addresses. On one side, heuristic methods such as greedy placement, threshold-based deferral, and rule-based routing dominate deployed systems because they are fast and simple to operate \citep{radovanovic2023, acun2023carbonexplorer, wijayawardana2025scheduling}. A growing body of work strengthens these heuristics with formal guarantees, as online and learning-augmented algorithms provide competitive bounds under uncertain inputs \citep{lechowicz2025stclip, bostandoost2024lacs}, meta-algorithms select among heuristics adaptively \citep{bostandoost2025datasched}, and distributionally robust or game-theoretic formulations offer probabilistic performance guarantees across a network of data centers \citep{hall2024probabilistic, breukelman2024hierarchical}.

On the other side, optimization-based approaches solve directly for a provably optimal schedule. Convex and linear programs were used early on to co-optimize workload and cooling \citep{liu2012} and, more recently, to match clean energy across space and time \citep{riepin2024spatiotemporal} and to minimize cost and carbon jointly across data-center micro-grids \citep{hao2024joint}. An exact optimum is valuable in this setting for two distinct reasons. It provides an upper bound against which any heuristic can be measured, and it guarantees that operational constraints are satisfied by construction. Bridging the two traditions, holistic design frameworks such as Carbon Explorer embed greedy scheduling within a larger exploration of operational and embodied trade-offs \citep{acun2023carbonexplorer}, while open testbeds increasingly support empirical comparison of the competing methods \citep{wiesner2024vessim}.

\subsection{Positioning of This Thesis}
\label{sec:lr-position}

Taken together, the literature reviewed above leaves a specific gap that this thesis addresses. Most deployed carbon-aware schedulers are heuristic and therefore cannot certify how close they come to optimal \citep{radovanovic2023, acun2023carbonexplorer}; most schedulers that exploit both temporal and spatial flexibility do so empirically or treat the two dimensions in isolation \citep{sukprasert2024limitations, asadov2025review}; and the closest exact joint formulation optimizes for dedicated 24/7 CFE procurement using synthetic generation profiles rather than dispatch against real grid signals \citep{riepin2024spatiotemporal}. In response to these three observations, this thesis formalizes carbon-aware scheduling as a deterministic linear program that is provably optimal by construction, optimizes the temporal and spatial dimensions jointly, draws on real two-year hourly CI and CFE data across five structurally diverse grid regions, and decomposes the realized saving to attribute it to spatial versus temporal flexibility. The resulting model is intended not as a replacement for the operational heuristics that dominate practice, but as the optimal reference they have so far lacked and as an instrument for understanding which form of flexibility actually drives carbon savings.
